\begin{center}\bfseries Resumen\end{center}
\begingroup
\fontsize{9.3}{10.4}\selectfont
\setlength{\parskip}{1.5pt plus .5pt minus .5pt}
\noindent
Se construye un operador positivo de masa sobre las fibras de partículas
persistentes y se determina su especialización escalar para el electrón
central. El resultado enlaza la identidad de la partícula, la clasificación
de sus estados y la evaluación de su masa: la magnitud procede de una
estructura que conserva la ruta, la orientación y la memoria de su
generación. En el sector electrónico, dos lectores bidireccionales
coinciden sobre una trayectoria explícita y fijan la misma corrección de
retorno; las demás colecciones conservan el alcance operatorio, espectral o
escalar demostrado para sus propias fibras.

La trayectoria central recorre exactamente tres diagonales nonádicas y
veintisiete celdas. Dos historias con el mismo histograma de valores absolutos
conservan memorias reversibles distintas. La transformada finita de la palabra
de retorno localiza tres modos no nulos y demuestra espectralmente
$(\Omega,P_6)=(80,6)$. Sobre las multisecciones, la ley relativa de masas
proporciona un criterio operatorio necesario y suficiente: un transporte TPK
biyectivo conserva la masa exactamente cuando conmuta con el operador de masa,
o, equivalentemente, cuando conserva su coordenada logarítmica completa.

El censo regional recorre \(104\,976\) pares de semillas, produce \(468\)
emisiones y retiene \(243\) palabras distintas. La biyección trítica por pares
sitúa esas palabras en el cuerpo \(S_0=E_9^3\) de \(729\) posiciones. La
diferencia exacta de umbral determina entonces la corona
\(243+27+1=271\) y la completación
\(729+243+27+1=1000\), con una normalización compatible bajo las proyecciones
entre dimensiones. La selección regional actúa después: ni la biyección ni
la completación cúbica definen una especie o un valor metrológico.

El marco propuesto es la Holografía Modular Triádica, en adelante HMT, y
su realización mediante Mecánica Dimensional. Su núcleo reúne tres objetos:
\emph{All Possible Paths}, en adelante APP, es una estructura aritmética
de caminos sobre \((\mathbb Z/9\mathbb Z)^2\): su grafo de Cayley aditivo,
con generadores cardinales, fija la adyacencia, y su realización celular
es toroidal. Las operaciones y las evaluaciones de suma y producto se
distinguen dentro de esa estructura;
TRIT es la unidad de firma orientada \(\{-1,0,+1\}\), que distingue los
regímenes locales; el \emph{Topological Phase Kernel}, en adelante TPK,
compone la dinámica de selección, transporte y memoria. Su evolución
enriquecida conserva hojas, residuos, cocientes, frontera y acarreo, y
produce la estructura discreta conjunta del continuo. La composición de
\(108\) actualizaciones determina el registro dodecafásico, y la transformada
orbital reconstruye reversiblemente \(K\), antecedentes de las coordenadas utilizadas por el lector
de masa.

La partícula se define por la persistencia de extensiones compatibles;
la ruta registra su historia observable. El atlas distingue \(81\) celdas,
\(56\) configuraciones de ruta, \(13\) multisecciones y \(324\) condiciones
orientadas. Se construyen sus cocientes, los observables de sabor y color,
la sección central única y las operaciones de conjugación, helicidad y
holonomía, incluida la banda de Möbius con su fibra y su lazo. La firma
integral y los momentos de los canales determinan el carácter positivo;
la acción, el reloj y la coordenatización de unidades proporcionan la sección
absoluta. Las razones de masa conservan el carácter relativo y las
diferencias entre exponentes de retorno.

La incidencia longitudinal marcada fija primero
\(N=5760\), \(r_N=5759/23040\) y
\(L_\alpha=\alpha_{\rm HMT}^{16}r_NL_\star\). De esa longitud se inducen
el reloj, \(R_{108}=L_\alpha\), la energía \(Q_L\), la masa cronológica y
la cotransformación del factor electrónico. Sobre cada fibra masiva, las
identidades \(H\Lambda=\hbar c_{\rm int}I\) y
\(\mathcal R_g\Lambda=L_\alpha^2I\) determinan de manera única
\(G=c_{\rm int}^3L_\alpha^2/\hbar\); su expresión circular es un
corolario posterior. La deducción no usa \(G\) como antecedente del
operador de masa ni modifica las tablas del espectro.

La composición se realiza sobre registros antes de evaluar sus firmas.
El desarrollo comprende sectores de masa nula, propagación orientada,
mezclas, resonancias y realizaciones orbitales de Kepler--Sommerfeld. La
completación graduada produce las reglas de ocupación y la exclusión de
Pauli sobre modos completos, con sus etiquetas de espín, hoja, sabor y
color. El contraste metrológico reúne cada evaluación con su ruta,
operador y coordenatización, y distingue valores centrales, intervalos, límites y
anchuras. Así se relacionan estructura, ecuación y valor conservando el
contenido específico de cada resultado y de su identificación física.

\noindent\textbf{Palabras clave:} derivación del espectro de masas;
estructura de las partículas; Holografía Modular Triádica; Mecánica
Dimensional; persistencia; holonomía; conjugación; comparación metrológica.
\par\endgroup
